\section{Conclusion}

Instant delivery in Indonesia will be won by whoever combines proximity, honesty of promise, and category ownership. This report specifies ParaExpress as ParagonCorp's answer: a two-speed network that preserves the NDC/DC backbone and superimposes a hyperlocal fulfillment tier built first from the consignment estate---a zero-capex coverage endowment formalized by Corollary~\ref{cor:stores}---and completed by cost-engineered dark MFCs sized by Theorem~\ref{thm:nodecount} and Proposition~\ref{prop:radius}.

The mathematics is not decoration. NP-hardness motivates the twin-gated design process; submodularity yields a provably near-optimal rollout order where capacity does not bind, and its failure under capacity (Section~\ref{sec:structural}-\ref{sec:capacitated}) is why the order is recomputed after every activation rather than fixed once; the uniform-matroid result makes weekly assortment refresh exact and cheap; total unimodularity guarantees that real-time routing is exactly solvable inside the sub-100\,ms wave budget; and the square-root laws of pooling and staffing---restated here under demand correlation and wave batching---quantify the density--decentralization trade that sank the first quick-commerce wave.

Revision~5 changes the report in one direction more than any other: it separates what is proved from what is assumed. Two results were being used outside the conditions under which they were proved, and both errors ran the same way, toward a network smaller and cheaper than the promise requires. One quantity that governs the headline promise---order-level availability---had no equation anywhere in Revision~4. The economics rested on a mean where a distribution was needed, and the acceptance gates could be passed by a model that was internally consistent and externally wrong. Correcting these does not weaken the thesis; it raises the capital ask and makes the thesis testable, which is the only form in which it is worth funding.

The competitive thesis inverts the platform playbook: rather than a platform using speed to commoditize brands, the brand uses speed---delivered from within the consumer's own neighborhood---to make the platform the commodity. Whether Indonesian beauty consumers will pay for that speed at the density this design requires is the one thing this report does not know, and experiment E7 exists to find out before the capital does.
